\chapter*{Résumé}
\addcontentsline{toc}{chapter}{Résumé}

Ce mémoire présente la conception et le développement de \projectname{}, une plateforme de recrutement orientée entreprise couvrant l'ensemble du cycle de recrutement : authentification, vivier de CV, gestion des offres, affectation des candidats, historique auditable des transitions, entretiens, analytique, administration, gouvernance des actions IA et onboarding. L'application repose sur une architecture \textit{feature-driven} construite avec Next.js 16, TypeScript, Bun, Better Auth, Neon PostgreSQL et Drizzle ORM.

L'apport majeur du projet réside dans sa couche d'intelligence artificielle, qui intervient à plusieurs niveaux : extraction des données CV, matching, screening, génération de guides d'entretien, agent conversationnel multi-outils, confirmation explicite des actions agentiques modifiant les données et pipeline RAG hybride avec citations. Cette couche conversationnelle suit désormais un modèle hybride, avec un assistant flottant contextuel et un workspace dédié \sourcecode{/agent} capable d'afficher des réponses sourcées, des traces d'exécution, des demandes de confirmation et des liens profonds vers les écrans métier. Les résultats d'évaluation réalisés dans le cadre du projet montrent une amélioration mesurable de la pertinence de recherche sur 40 requêtes, avec une couverture de vérité terrain de 85\%.
L'expérience agentique finale ajoute aussi une restitution visuelle des analyses : lorsqu'une question porte sur une tendance, un pipeline, des compétences, la demande de postes ou un écart demande--offre, la réponse peut afficher des cartes graphiques construites à partir des mêmes sorties d'outils vérifiées, au lieu de se limiter à un texte descriptif.

Le projet démontre qu'une IA utile au recrutement doit rester ancrée dans des données structurées, des workflows explicites, des historiques immuables et des garde-fous techniques. La plateforme obtenue constitue ainsi une base crédible pour une digitalisation avancée et gouvernée des opérations de recrutement.
